\begin{textsample}{POS Dim 7 – global\_south\_2025\_02 – Score 3.00 – global\_south\_2025\_02/120893817.txt}  \label{ex:f7_pos_423}
President Donald Trump ' s proposal for the U.S. to take over Gaza and resettle Palestinians has been rejected globally.
The plan, aimed at transforming Gaza into an economic hub, faced strong opposition from international leaders, raising concerns over worsening the Israeli-Palestinian conflict.
APUS President Donald Trump wants the US to ' take over ' Gaza.
President Donald Trump ' s proposal for the United States to"take over"the Gaza Strip and resettle its Palestinian residents was quickly met with widespread rejection and condemnation from around the world on Wednesday.
During a White House news conference with Israeli Prime Minister Benjamin Netanyahu, Trump unveiled a plan that he described as transforming Gaza into"the Riviera of the Middle East."He announced that the U.S. would take"ownership"of Gaza, dismantle dangerous unexploded bombs, level the war-torn territory, and create an economic hub to provide"unlimited numbers of jobs.""The U.S. will take over the Gaza Strip, and we will do a job with @ @ @ @ @ @ @ @ @ @ would be responsible for rebuilding Gaza and overseeing its redevelopment.
However, the proposal sparked strong opposition from various nations, stirring concerns that it would complicate ongoing ceasefire talks and deepen the Israeli-Palestinian conflict.
Trump ' s proposal : What the world said Turkey : Turkey swiftly condemned Trump ' s comments.
Turkish Foreign Minister Hakan Fidan called the U.S. proposal"unacceptable,"warning that any plan that disregarded Palestinian interests would only escalate the conflict.
Fidan told Anadolu news agency that Turkey would \textbf{reconsider} its actions against Israel, including cutting off trade and recalling its ambassador, if Israel ceased its killing of Palestinians and improved their conditions.
China : China also voiced its opposition to Trump ' s plan.
A Chinese Foreign Ministry spokesperson stated that Beijing opposes the forced relocation of people in Gaza.
Emphasizing China ' s longstanding support for a two-state solution, the spokesperson reiterated that Palestinian governance over Gaza remains a crucial principle for the future of the region."Palestinian rule is the basic principle of post-war governance in Gaza @ @ @ @ @ @ @ @ @ @ for a peaceful resolution that respects Palestinian sovereignty.
Houthi Rebels : The leader of Yemen ' s Houthi rebels, Mohammed al-Bukhaiti, criticized Trump ' s remarks as an expression of"American arrogance."On X, he expressed that if Egypt or Jordan decided to oppose the U.S., Yemen would stand in full support against the plan.
The Houthis, who have previously targeted Israel and shipping through the Red Sea, warned that Trump ' s proposal could lead to greater regional instability if met with resistance.
U.S.
House Speaker ' s Support : On the domestic front, U.S.
House Speaker Mike Johnson, a \textbf{Republican} from Louisiana, praised Trump ' s approach, calling it a"bold action"aimed at bringing"lasting peace"to Gaza.
He expressed hope that the plan would stabilize the region, demonstrating domestic support for Trump ' s ambitious proposal.
Secretary of State Marco Rubio : U.S.
Secretary of State Marco Rubio also expressed his support for Trump ' s comments.
Rubio wrote on X,"Gaza MUST BE @ @ @ @ @ @ @ @ @ @ a revitalized Gaza, jokingly referring to it as part of an effort to"Make Gaza Beautiful Again,"a reference to Trump ' s campaign slogan.
Rubio stressed the U.S. ' s commitment to peace and rebuilding the region.
Middle Eastern and Global Criticism Despite the backing from some U.S. officials, the proposal was met with fierce criticism from many Middle Eastern countries, particularly Saudi Arabia and others in the region.
These nations, which have long advocated for a Palestinian state in both Gaza and the West Bank with East Jerusalem as its capital, argued that Trump ' s plan undermines Palestinian sovereignty.
The global consensus remains largely in \textbf{favor} of a two-state solution that would allow Palestinians self-determination, and the U.S. proposal to take over Gaza was seen as a violation of that principle.

% matched lemmas: favor, reconsider, republican
\end{textsample}
